\chapter{実装}

本章では，第 3 章で述べた設計に基づいて構築した評価基盤の実装について述べる．4.1 節で実行環境を，4.2 節で全体の構成を示し，4.3 節から 4.6 節で各モジュールについて述べる．4.7 節では実装の過程で判明した不具合とその対処を述べ，4.8 節で本章をまとめる．

\section{実行環境}

実行環境を表 4.1 に示す．

\begin{table}[htbp]
\centering
\small
\caption{実行環境}
\label{tab:4.1}
\begin{tabularx}{\linewidth}{|l|X|}
\hline
項目 & 内容 \\
\hline
OS & Windows 11 \\
言語 & Python 3.12 \\
主要ライブラリ & anthropic，numpy，matplotlib \\
モデルへの接続 & Anthropic Messages API \\
バージョン管理 & Git / GitHub \\
\hline
\end{tabularx}
\end{table}

外部ライブラリへの依存は意図的に小さくした．較正指標の計算は scikit-learnなどを用いず numpy のみで実装している．指標の定義を式から直接実装することで，ライブラリの仕様に依存した解釈の揺れを避けるためである．

\section{システム構成}

実装はモジュールごとにファイルを分けた．構成を表 4.2 に示す．

\begin{table}[htbp]
\centering
\small
\caption{モジュール構成}
\label{tab:4.2}
\begin{tabularx}{\linewidth}{|X|X|}
\hline
ファイル & 役割 \\
\hline
\texttt{fetch\_\allowbreak{}mgsm.\allowbreak{}py} & データセットの取得と実験用形式への変換 \\
\texttt{prompts.\allowbreak{}py} & 引き出し方式ごとのプロンプト生成 \\
\texttt{run\_\allowbreak{}pilot.\allowbreak{}py} & モデルへの問い合わせと結果の記録 \\
\texttt{parser.\allowbreak{}py} & 応答の解析と正誤判定 \\
\texttt{calibration.\allowbreak{}py} & 較正指標と Murphy 分解の計算 \\
\texttt{compare\_\allowbreak{}methods.\allowbreak{}py} & 方式間の比較と作図 \\
\texttt{bootstrap\_\allowbreak{}ci.\allowbreak{}py} & 信頼区間の算出 \\
\texttt{rescore.\allowbreak{}py} & 保存済み応答からの再採点 \\
\texttt{test\_\allowbreak{}parser.\allowbreak{}py}，\texttt{test\_\allowbreak{}calibration.\allowbreak{}py} & 単体テスト \\
\hline
\end{tabularx}
\end{table}

この分割により，解析の不具合を修正した際に，モデルへの問い合わせをやり直さずに結果を再計算できる．実際にこの仕組みが必要になった経緯は 4.7 節で述べる．

\section{プロンプト生成}

\texttt{prompts.\allowbreak{}py} は，引き出し方式と問題形式の組み合わせに応じてテンプレートを選び，問題文を埋め込む．問題形式は自由回答（数値），多肢選択，翻訳の 3 種を想定し，それぞれに専用のテンプレートを用意した．

多肢選択のテンプレートは「選択肢の中から選び，A から D の記号で答えよ」と指示する．ここで，選択肢が空のまま多肢選択のテンプレートを使うと，選ぶ対象が提示されないままモデルが記号を答えることになる．この状態はエラーにならず，生成された結果が一見正常に見えるため発見が遅れる．そこで，多肢選択の問題で選択肢が空の場合は明示的に例外を送出する検査を実装した．

Verb.2S については，第 1 ターンと第 2 ターンの 2 つのプロンプトを返す．第 2 ターンのプロンプトには，第 1 ターンで得られた回答を文字列として埋め込む．

\section{モデルへの問い合わせ}

\texttt{run\_\allowbreak{}pilot.\allowbreak{}py} は問題を 1 問ずつ処理し，応答を CSV に記録する．記録する項目は，問題の識別子，領域，モデル名，方式，問題文，選択肢，回答形式，正解，モデルの回答，確信度の生の値，正規化した確信度，正誤，解析の成否，そしてモデルの応答全文である．

応答全文を保存しているのは，解析の不具合が後から見つかった場合に，API を呼び直さずに再採点できるようにするためである．

ただし Verb.2S については保存が不完全である．記録される応答全文は第 2 ターンのものであり，第 1 ターンの応答は抽出後の回答だけが残る．したがって Verb.2Sでは回答の再抽出ができず，回答の抽出に関わる不具合を後から修正できない．この制約は 4.7 節で述べる不具合の検証範囲に影響した．今後は両ターンの応答全文を別の欄に保存する．

一時的な通信エラーやレート制限に対しては，待ち時間を倍にしながら再試行する．ただし型の誤りや設定の不備といった，再試行しても解消しない種類の例外は再試行の対象から除外した．これらを再試行すると，即座に判明するはずの誤りの発見に数十秒を要することになるためである．

温度の指定については，モデルの世代によって扱いが異なる．サンプリングの指定を受け付けるモデルにのみ温度を送り，受け付けないモデルには送らない．どちらとも判断できないモデル名が指定された場合は，温度を送らずに警告を表示する．温度を指定しない場合は出力に揺らぎが生じ，予備実験が前提としている条件から外れるため，その事実を実行者に伝える必要がある．

\section{応答の解析と正誤判定}

\texttt{parser.\allowbreak{}py} は，モデルの応答から回答と確信度を抽出する．モデルが指定された形式から外れた応答を返すことがあるため，複数の正規表現を順に試す実装とした．また，モデルが仮の記入をしてから途中式を書き，最後に本来の答えを書く場合があるため，回答は応答中で最後に現れたものを採る．このとき，見出しの種類（「回答:」「Answer:」など）の優先順ではなく，応答中の出現位置で判断する．見出しが混在する応答で，前方にある見出しの側を拾わないようにするためである．見出しの区切りには半角と全角の両方のコロンを認める．

確信度に百分率記号が付いている場合は，値域によらず 100 で除す．値域だけで尺度を判断すると，「確信度: 0.5\%」を 0.5，すなわち 50\% と読み違える．

確信度の解析は方式によって異なる．Verb.1S と Verb.2S では数値を抽出する．プロンプトでは 0.0 から 1.0 の範囲を指示しているが，モデルが 0 から 100 の尺度で答える場合があるため，値が 1.0 を超えるときは 100 で除して正規化する．Ling.1S では表 3.1 の言語表現を抽出し，対応する数値に変換する．このとき，短い表現が長い表現の一部に一致して取り違えが起きないよう，文字数の多い表現から順に照合する．

正誤判定は回答形式に応じて切り替える．数値の場合は，**回答の数値部分どうしを比較する簡略な採点規則**を用いる．正解と回答の双方に同一の正規化を施し，正規化後の文字列全体が単一の数値として読める場合にかぎり数値として比較する．許容する表記は以下に限る．

\begin{itemize}
\item 全角の数字と記号（半角に変換する）
\item 囲みの括弧 1 組（\texttt{[32]} など）
\item 末尾の百分率記号 1 つ
\item 桁区切りのカンマ（3 桁区切りとして正しい場合にかぎり除去する）
\end{itemize}

正規化後に数値として読めない回答は受理しない．分数（\texttt{1/\allowbreak{}2}），指数表記（\texttt{1e3}），小数点前の 0 を省いた表記（\texttt{.\allowbreak{}5}），複数の値を含む回答（\texttt{32 or 64}）はいずれも未対応であり，誤答として扱う．また，正解を数値として読めない場合はデータ側の不備として例外を送出する．文字列の一部に数字があれば採用する方式は採らない．その方式では \texttt{1/\allowbreak{}2} を \texttt{1} として受理してしまい，表記の正規化を超えて意味の異なる回答まで正答としてしまうためである．

\textbf{この規則は単位の整合性を保証しない．} 単位の語（\texttt{円}，\texttt{個}，\texttt{cm} など）は取り除かない．正解欄に単位が無いことは，回答側の単位を消してよい理由にはならないためである．たとえば 100 ドルを求める設問で「100セント」と答えた場合，単位を消すと正解と一致してしまう．単位を伴う回答は，設問が要求する単位と照合しなければ正誤を判定できず，本研究の採点規則はそこまで踏み込んでいない．

百分率記号のみ例外としている．予備実験で用いた 250 問のうち，百分率を伴う回答が現れたのは確率を問う 1 問（mgsm\_ja\_227）であり，正解が単位のない数値 33 で与えられている．この設問を確認したうえで，\texttt{33\%} を同じ値の別表記として扱った．一般の単位に拡張できる判断ではない．

なお予備実験のデータでは，\textbf{単位を伴う回答は 1 件も現れなかった}．数値そのものでない回答は 750 件中 9 件で，内訳は角括弧を伴うもの 7 件と百分率記号を伴うもの 2 件である．したがって上記の制約が予備実験の結果に影響した箇所はない．

多肢選択の場合は記号を選択肢の本文に解決してから照合する．データセットによって正解が記号で与えられる場合と選択肢の本文で与えられる場合があるため，両方に対応している．なお現在の実装は選択肢を 4 つまで想定しており，5 つ以上の選択肢を持つデータセットを用いる際は記号の生成規則を改める必要がある．

\section{指標の計算}

\texttt{calibration.\allowbreak{}py} は 3.4 節の定義をそのまま実装している．ECE は確信度を10 区間に等分し，各区間の平均確信度と正答率の差を標本数で重みづけて合計する．Murphy 分解も同じ区間分割を用いる．

\texttt{bootstrap\_\allowbreak{}ci.\allowbreak{}py} は信頼区間を算出する．問題を復元抽出して指標を計算し直す操作を繰り返し，得られた分布の 2.5 パーセンタイルと 97.5 パーセンタイルを区間の端とする．

ここで重要なのは，3 方式がまったく同じ問題集合に対して実行されている点である．したがって再抽出は方式ごとに独立に行わず，問題ごとに 3 方式の結果を組にして同じ添字で抽出する．こうすることで，問題単位の対応関係を保ったまま差の分布を求められる．方式ごとに独立に抽出すると，この対応関係が失われ，差の分散が方式間の共分散を反映しなくなる．

AUROC は，再抽出の結果すべてが正答またはすべてが誤答になった場合に計算できない．この場合は当該の反復を除外し，有効な反復数を併記する．少数の標本から求めた区間を，そうでないものと取り違えないようにするためである．

\section{実装過程で判明した不具合}

予備実験の実施にあたり，結果を無効にする不具合が 6 件見つかった．いずれも実行時にエラーを発生させず，出力された数値が一見正常に見えるものであった．再発を防ぐため記録しておく．前半の 4 件は実験の実施前に，後半の 2 件は外部のレビューを受けて実施後に見つかったものである．

\textbf{選択肢を持たない問題への多肢選択テンプレートの適用．} 自作した予備データにおいて，多肢選択の領域に分類されているにもかかわらず選択肢が空の問題が100 問中 42 問存在した．これらには選択肢が提示されないまま「記号で答えよ」という指示のみが送られ，モデルはすべて同じ記号を返していた．正答率はその記号が正解である問題の割合に一致し，モデルの能力とは無関係な値になっていた．4.3 節で述べた検査を追加し，このデータの使用を取りやめて MGSM に変更した．

\textbf{確信度の見出し語の不一致．} プロンプトは「確信度」という語で確信度を求めているのに対し，解析側は「信頼度」「自信度」のみを照合していた．その結果，確信度が 1 問も抽出できていなかった．

\textbf{尺度の二重変換．} 確信度を 0 から 100 の尺度として一律に 100 で除していたため，0.9 と答えられた場合に 0.009 として扱われていた．値域から尺度を判定する実装に改めた．

\textbf{Ling.1S における確信度の欠損．} 方式によらず数値用の解析を適用していたため，言語表現で返される Ling.1S の確信度が 1 件も抽出できていなかった．方式ごとに解析方法を切り替える実装に改めた．応答全文を保存していたため，API を呼び直すことなく \texttt{rescore.\allowbreak{}py} で再採点できた．

\textbf{正解側の桁区切りカンマを除去していなかった．} 数値の比較において，回答側だけカンマを除去し，正解側はそのまま数値へ変換していた．正解が \texttt{2,125} の形式であると変換に失敗し，例外を捕捉して部分一致の判定に移る．\texttt{2,125} は\texttt{2125} に含まれないため，正しい回答が誤答として記録されていた．3 方式を通じて 4 件確認された．

\textbf{応答中で最初に現れた回答を採っていた．} モデルが「回答: [計算します]」のように仮の記入をしてから途中式を書き，最後に「回答: 64」と本来の答えを書く場合，仮の記入のほうが回答として記録されていた．確認できた範囲で Verb.1S に2 件，Ling.1S に 13 件発生しており，\textbf{方式によって発生数が偏っていた}．そのため修正前の正答率の差には，採点処理の影響が混入していた．なお 4.4 節で述べたとおり Verb.2S は第 1 ターンの応答を保存していないため，同種の不具合が何件あったかを確認できない．

後半の 2 件は，\texttt{rescore.\allowbreak{}py} により保存済みの応答から再採点した．再採点の前後で正答数は Verb.1S が 211 から 215，Verb.2S が 218 から 221，Ling.1S が181 から 195 に変化した．第 5 章に示すのはすべて修正後の値である．また，同種の誤りを検出するための回帰テストを \texttt{test\_\allowbreak{}parser.\allowbreak{}py} に追加した．

これらはいずれも，エラーを出さずに誤った数値を出力する種類の不具合である．指標が計算できたことをもって処理が正しいと判断してはならない．中間的な出力，とくにモデルの応答全文と解析結果を目視で確認する工程が必要である．また，外部のレビューによって 2 件が新たに見つかったことは，実装者自身による確認だけでは不十分であることを示している．

\section{本章のまとめ}

本章では評価基盤の実装について述べた．モジュールを分割し，モデルの応答全文を保存する設計としたことで，解析の不具合が後から見つかった場合にも API を呼び直さずに再計算できる．実際に 2 件の採点の不具合を，追加の問い合わせなしに修正できた．ただし Verb.2S については第 1 ターンの応答を保存しておらず，この利点が得られていない．また，判明した 6 件の不具合がいずれもエラーを伴わずに誤った数値を生じさせるものであったことを記録した．次章では，この基盤を用いた予備実験の結果を示す．
